\chapter{La presencia material del Estado, 1923--1945}
\label{cap:presencia}

Otros capítulos de este libro, anteriores y posteriores, muestran al Estado provincial decidiendo
sobre el departamento: aprueba presupuestos sin publicarlos, concede agua sin medirla,
categoriza municipios y retira caminos de un plan (capítulos \ref{cap:hacienda} y
\ref{cap:infraestructura}). Este capítulo pregunta otra cosa, más simple y más difícil de contestar:
\textbf{qué había ahí}.

No qué se decidía sobre el territorio, sino qué edificios, qué cargos, qué dotaciones y qué
servicios tenía el departamento, y cuánto costaban. La respuesta sale casi entera del barrido
continuo del Boletín entre 1923 y 1945, y tiene la ventaja de que \textbf{en ese tramo lo que no
aparece es un resultado y no un hueco de búsqueda} (apéndice \ref{cap:fuentes}).

\begin{cautela}
Una advertencia sobre el método, porque cambia lo que las cifras significan. El Boletín publica
\textbf{actos de gasto}, no inventarios: una partida liquidada prueba que se autorizó un pago,
no que la cosa llegara ni que siguiera funcionando. Este capítulo cuenta lo que el Estado
decidió poner, y en los pocos casos en que el archivo permite verlo, también si llegó.
\end{cautela}

\section{La escuela}

En 1934 la Ley 146 de presupuesto de educación registra, entre las «Escuelas Elementales de 2\textsuperscript{a}
Categoría», la \textbf{Escuela «Juana Moro de López» de La Caldera}, con una planta de tres
personas: directora a \$125 y dos maestras a \$100 cada una. \textbf{\$325 mensuales, \$3.900 al
año}, para todo el pueblo. En Vaqueros, la escuela «Bernabé López». Son las dos escuelas
provinciales del departamento, y siguen siendo las dos en el presupuesto de 1935.

Puesto contra la hacienda municipal, ese número ordena la escala: la escuela provincial del
pueblo cuesta al año más de una vez y un tercio lo que el capítulo \ref{cap:hacienda} infiere como renta
anual del municipio en 1937.

\textbf{Todo lo demás en materia escolar lo puso la Nación.} Entre 1937 y 1945 el archivo
registra cuatro trámites de escuelas de la \textbf{Ley 4874, la Ley Láinez}, por las cuales el
Consejo Nacional de Educación instalaba escuelas en provincias con aquiescencia provincial:

\begin{itemize}[leftmargin=*,itemsep=2pt]
\item \textbf{Potrero de Castilla, 1937.} La Provincia niega la aquiescencia alegando que
instalará escuelas fiscales propias; \textbf{veinticuatro días después} deroga ese artículo y la
concede, porque la Inspección Nacional verificó el analfabetismo y comprobó que \textbf{no había
escuela provincial en la localidad ni en lugares próximos}.
\item \textbf{Gallinato, 1939.} Decreto 3456: aquiescencia concedida.
\item \textbf{Yacones, 1940.} Decreto 90: aquiescencia concedida.
\item \textbf{San Alejo, octubre de 1945.} Decreto 9156-G: \textbf{51 niños en edad escolar y la
escuela más cercana a diez kilómetros}.
\end{itemize}

El mismo día de octubre de 1945, el Decreto 9155-G autoriza \textbf{trasladar la Escuela Nacional
N\textsuperscript{o} 250 de Wierna} a la localidad de San Francisco, partido de Yacones.

\textbf{Leídos juntos, los cuatro trámites dicen una sola cosa}: en los parajes del departamento
---no en el pueblo--- la escuela la puso la Nación, y en el único caso en que la Provincia dijo
que la pondría ella, no era cierto y hubo que corregirlo en menos de un mes.

\textbf{La escuela del pueblo, en cambio, es anterior a este capítulo.} En julio de 1911 el
Ministerio de Hacienda (B.O.\ Nº 268, 1911, h.\ 3) jubila a \textbf{doña Jacoba Luquez de Gómez, directora de la «Escuela
elemental de la parroquia de la Caldera»}, con \$78,74 mensuales calculados sobre el promedio de
sus últimos cinco años. La jubilación prueba dos cosas: que en
1911 había escuela en el pueblo y que su directora llevaba allí, por lo menos, un lustro. La
«Juana Moro de López» del presupuesto de 1934 no es, entonces, la primera.

\section{La comisaría, y sus subcomisarías}

Es la institución mejor documentada del período, por lejos: el archivo registra decenas de
nombramientos, renuncias, licencias, permutas, reconocimientos de servicios y planillas de
personal supernumerario. El departamento tuvo comisaría en el pueblo y subcomisarías en
\textbf{Vaqueros} ---ad honorem hasta que en 1936 pasa a rentada---, \textbf{Mojotoro},
\textbf{San Alejo y Santa Rufina} y, desde 1934, \textbf{Lesser}.

Tres datos le ponen escala.

\textbf{El teléfono.} Una factura de la Compañía Argentina de Teléfonos de diciembre de 1933
consigna «\textit{Telef.\ N\textsuperscript{o} \textbf{7} --- Policía de la Caldera \dots{} 2,40}», dentro de un total
reconocido de \$59,55. Un número de un solo dígito, contra números de cuatro en la capital: la
policía del departamento estaba conectada, y era de las primeras líneas de la red.

\textbf{El edificio.} En 1944 la refección del local de la comisaría del pueblo se licita, se
adjudica a único oferente y se certifica en ocho actos, por \$5.573,80 finales
(capítulo \ref{cap:defensas}). Los expedientes más antiguos del legajo son de 1942: \textbf{la
refacción vino tramitándose por lo menos dos años}.

\textbf{Y el terreno.} En 1943 el Estado tuvo que promover \textbf{juicio de posesión
treintañal} para adquirir los ocho mil metros cuadrados donde su comisaría funcionaba desde
1913 o antes. Ese edificio tiene un antecedente anterior al período de este capítulo: \textbf{el
15 de noviembre de 1909 el Gobierno de la Provincia llamó a licitación «para la construcción de
un edificio destinado á la Policía de La Caldera»}, con planos y pliego en la Oficina de
Inspección de Obras Públicas y apertura el 25 de ese mes (B.O.\ Nº 107, 1909). Nueve meses después, el 1º de agosto de 1910, un decreto aprueba el \textbf{contrato ad
referendum} celebrado entre el jefe de la Oficina de Obras Públicas y \textbf{don Moisés Vera}
«\textit{sobre construcción del edificio para la comisaría de La Caldera}» (B.O.\ Nº 178,
1910). El decreto no dice monto, ni plazo, ni en qué lugar del departamento, y el contrato no se
publica. El archivo no dice qué se construyó, y el llamado es de los mismos años en que se licitan
comisarías en Cerrillos, Coronel Moldes, La Viña y La Merced: la comisaría de La Caldera no fue
una excepción sino parte de un plan. Si lo contratado en 1910 es lo que estaba en pie en 1913, el
Estado lo construyó sobre un terreno que recién treinta y tres años después intentó poner a su
nombre. El expediente de esa licitación está pedido (apéndice \ref{cap:pedidos}). \textbf{Y el juicio de 1943 le costó al Estado ochenta y dos pesos}: setenta al diario por el edicto y doce al Juez de
Paz por diligenciar un oficio. \textbf{El único edificio público del que el archivo permite
seguir la historia completa es uno que el Estado ocupaba sin título desde hacía treinta años.}

\section{Los dos escalones más bajos}

En marzo de 1934, en cinco días, dos decretos clasifican las oficinas del Estado en el
departamento.

El \textbf{Decreto 17606} fija las categorías de las oficinas del Registro Civil: La Caldera,
\textbf{tercera de tres}. El \textbf{Decreto 17640} hace lo mismo con las Receptorías de Rentas:
La Caldera, \textbf{quinta de cinco}. Un año antes, la nómina de libros del Registro Civil
consigna la oficina «Caldera» con su libro de matrimonios.

\textbf{Y el primer sueldo policial que el archivo pone en el pueblo es de 1912.} En diciembre de
ese año el Poder Ejecutivo nombra \textbf{interinamente subcomisario de policía del pueblo de La
Caldera} a Silverio Cáseres (B.O.\ Nº 379, 1912, h.\ 3), «\textit{con el sueldo de cien pesos mensuales}». Es el único
nombramiento de la época que distingue el pueblo del departamento, y la cifra ordena la escala
como ninguna otra: \textbf{ese cargo costaba al año \$1.200, casi lo mismo que toda la renta
municipal del departamento en 1908} (capítulo \ref{cap:hacienda}). Un subcomisario interino valía
un municipio entero.

Esa oficina única tiene un antecedente que explica la geografía del departamento mejor que
cualquier mapa. En octubre de 1911 el Poder Ejecutivo creó \textbf{otra oficina del Registro
Civil en la estación Mojotoro} (B.O.\ Nº 291, 1911, h.\ 3), ad honorem y sólo para el enrolamiento, porque la separaba del
pueblo «\textit{la larga distancia}» y por «\textit{las dificultades con que tropieza la gente
para trasladarse a aquel pueblo por la falta de medios de movilidad y de transporte}». La puso a
cargo de \textbf{Benjamín López}\index{López, Benjamín}, que ese mismo año era el subcomisario
del partido de la estación. En diciembre de 1912 el mismo López pasa a \textbf{comisario del
partido de Mojotoro}, por renuncia de su antecesor (B.O.\ Nº 379, 1912, h.\ 3), y en el juicio
de ese año declara como \textbf{jefe de la estación} del ferrocarril: registro civil, policía y
estación, en la misma persona y en el mismo paraje. Un cargo sin sueldo, en la punta del departamento a la que llegaba el ferrocarril:
el Estado llegaba donde llegaba el tren, y gratis.

La Receptoría existió durante todo el período y su titular cobraba por comisión sobre lo
recaudado, no por sueldo: el archivo liquida esas comisiones en 1923, 1926, 1930, 1931, 1935,
1942 y 1943, y en varios años la Provincia tuvo que sancionar \textbf{leyes de crédito
suplementario} ---las Leyes 560, 586, 605 y 705--- para pagar lo que debía a la receptora o a la
expendedora de La Caldera. \textbf{Es una de las pocas series continuas del archivo, y lo que
registra es una deuda recurrente del Estado con quien le cobraba los impuestos.}

Con el Decreto 14205 de 1931, que clasifica al municipio en \textbf{tercera categoría}, y con la
ubicación del departamento en la \textbf{Primera Zona policial} de 1943 ---la única de siete sin
inspección propia ni asiento---, el cuadro se cierra: \textbf{en todas las escalas que el Estado
usó para ordenarse, La Caldera quedó abajo o afuera.}

\section{Lo que se votó y lo que llegó}

Cuatro actos del período prometen o entregan una cosa material al pueblo. Conviene leerlos en
orden y con sus montos.

{\small
\begin{longtable}{@{}p{1.4cm}p{5.5cm}p{2.2cm}p{3.5cm}@{}}
\toprule
\textbf{Año} & \textbf{Qué} & \textbf{Monto} & \textbf{Estado} \\
\midrule
\endhead
1929 & Terreno en el pueblo para la \textbf{Casa Municipal} (Ley 11088) & \textbf{\$500} & Autorizado \\
\addlinespace
1935 & \textbf{Iglesia Parroquial}: reforma y ampliación, planos del ing.\ Antonio Mille (Ley 172) & \textbf{\$7.000} & Autorizado \\
\addlinespace
1935 & \textbf{Estación sanitaria tipo «B»} en Caldera, dentro del plan de obras públicas por empréstito (Ley 158) & \textbf{\$18.360} & En el cuadro de la ley. \textbf{Sin rastro posterior} \\
\addlinespace
1937--42 & \textbf{Gimnasio completo} para el pueblo, por la Junta de Educación Física (Ley 439) & \textbf{\$1.500} & \textbf{Entregado}: lo recibe Augusto Regis\index{Regis, Augusto} \\
\bottomrule
\end{longtable}}

La tabla se lee sola, pero tres cosas merecen decirse.

\textbf{El gimnasio es el único de los cuatro cuya llegada está probada.} La Ley 439 de
septiembre de 1937 destinó \$55.000 a comprar y colocar gimnasios en pueblos de la provincia; en
enero de 1942 la Junta de Educación Física paga dos, a \$1.500 cada uno, provistos a La Silleta
y a La Caldera, y el expediente registra que \textbf{Augusto Regis, Presidente de la Comisión
Municipal, acusó recibo} por comunicaciones del 24 de diciembre de 1941 y del 7 de enero de
1942. Es, de los cuatro, el único caso en que el archivo permite decir \textit{llegó}; fuera de la tabla, la refacción de la comisaría de 1944, certificada en ocho actos, es el otro.

\textbf{La estación sanitaria es el único caso en que se puede decir lo contrario.} Figura en el
cuadro de inversiones de la Ley 158 de 1935 con monto asignado, y \textbf{no vuelve a aparecer
en las mil cuatrocientas setenta y dos ediciones del barrido}. Lo que sí aparece, en febrero de 1936, es un
decreto que provee al departamento \textbf{un guarda sanitario con asiento en La Caldera, por
treinta días, con asignación de \$150 mensuales}. Treinta días, una persona, \$150: eso es lo que
el archivo muestra de la salud pública del departamento en el período.

\textbf{Y los \$500 de la Casa Municipal de 1929 valen como unidad de medida.} Es lo mismo que
seis años después se presupuestará para el espigón de emergencia para que el río no entre a la calle
principal, y la catorceava parte de lo que la Legislatura votó el mismo año de ese espigón para
reformar la iglesia.

\section{Lo que sí se construyó sin pausa}

Hay una excepción grande a todo lo anterior, y es la que el capítulo \ref{cap:infraestructura}
desarrolla: \textbf{los caminos}.

Entre 1930 y 1942 el archivo sigue, acta por acta del Directorio de Vialidad, la construcción
del camino de \textbf{La Calderilla a El Desmonte y a Campo Santo por Gallinato} ---con puentes,
alcantarillas, un puente de veinte metros de luz, alambrado, portland, sobrestantes, cuadrillas
y muertos---, del de \textbf{Salta a Los Yacones por San Lorenzo y Castellanos}, del de
\textbf{Vaqueros a Lesser}, del de \textbf{Lesser a Los Yacones}, del de \textbf{Yacones a
Potrero de Castillo}, del de \textbf{La Caldera al Sauce} y del tramo del \textbf{camino nacional
de Salta a Jujuy} que pasa por el pueblo. En 1940, el plan de quince años de la Ley 12.625
asigna \$200.000 a un \textbf{Puente Río Caldera}.

\textbf{La diferencia con todo lo demás de este capítulo es de orden de magnitud.} La escuela
del pueblo costaba \$3.900 al año; el gimnasio, \$1.500; el guarda sanitario, \$150 al mes; el
puente proyectado, doscientos mil. \textbf{Lo que el Estado construyó en este departamento entre
1923 y 1945 fueron caminos}, y los caminos no terminan acá: van a Campo Santo, a Güemes, a
Jujuy, a la capital.

\section{La conclusión de este capítulo}

Puesto todo junto, el inventario de la presencia material del Estado en el departamento entre
1923 y 1945 es corto y tiene una forma reconocible.

\textbf{Hubo instituciones, y fueron las cuatro más baratas}: una comisaría con cuatro
subcomisarías, una oficina de Registro Civil de tercera categoría, una Receptoría de quinta y
un Juzgado de Paz lego con vacantes recurrentes. Ninguna requiere edificio propio; de hecho, el
único edificio del que hay constancia es el de la comisaría, y el Estado lo ocupaba sin título.

\textbf{Hubo una escuela provincial en el pueblo y otra en Vaqueros}, y en los parajes las
escuelas las puso la Nación.

\textbf{Hubo una promesa de estación sanitaria que no dejó rastro}, y un guarda sanitario por
treinta días.

\textbf{Y hubo caminos, muchos, caros y continuos.} Ahí se concentró la inversión.

Este capítulo no concluye que el Estado estuviera ausente: estuvo, de manera continua y
documentada, durante veintitrés años. \textbf{Concluye que estuvo de una forma determinada}
---policía, registro, recaudación y vías de paso--- y que las dos cosas que un pueblo necesita
para quedarse ---salud y escuela más allá del casco--- fueron, una, una partida sin ejecución
comprobable, y la otra, cosa de la Nación.

Es el mismo reparto que el capítulo \ref{cap:infraestructura} encuentra ochenta años después,
con otras palabras y otros montos.
